%% ma: L12-B14-0055
%% lop: 12
%% chuong: 5
%% bai: 14
%% dang: 12.14.1
%% loai: TN4
%% muc: NB
%% dap_an: "A"
%% nguon: "To_n_12_PT_M_t_Ph_ng.pdf (Đề kiểm tra đánh giá kiến thức, Chuyên đề Phương trình mặt phẳng, Nguyễn Anh Khôi - Lê Minh Khôi), Câu 9"
%% kien_thuc: "Bài 14 (phương trình tổng quát của mặt phẳng, vectơ pháp tuyến)"
%% trang_thai: chinh-thuc
%% ghi_chu: "Tài liệu gốc không kèm đáp án; đáp án do Claude giải và kiểm bằng sympy. Câu 35 của tài liệu gốc trùng hoàn toàn câu này (chỉ đảo thứ tự phương án), đã bỏ câu 35."
\begin{ex}
Trong không gian $Oxyz$, cho mặt phẳng $(P)\colon 2x-4y+7=0$. Chọn khẳng định đúng.
\choice{\True Mặt phẳng $(P)$ có vô số vectơ pháp tuyến và $\vec n_1=(2;-4;0)$ là một vectơ pháp tuyến của $(P)$.}{Mặt phẳng $(P)$ có duy nhất một vectơ pháp tuyến, vectơ đó là $\vec n_2=(2;-4;7)$.}{Mặt phẳng $(P)$ có duy nhất một vectơ pháp tuyến, vectơ đó là $\vec n_1=(2;-4;0)$.}{Mặt phẳng $(P)$ có vô số vectơ pháp tuyến, trong đó có một vectơ là $\vec n_2=(2;-4;7)$.}
\loigiai{$(P)\colon 2x-4y+0z+7=0$ có vectơ pháp tuyến $(2;-4;0)$; mọi vectơ $k(2;-4;0)$, $k\ne0$, cũng là vectơ pháp tuyến nên có vô số. Chọn A.}
\end{ex}
